\section{Results}

Taking logs of~\eqref{eq:production} and differencing over time gives the usual
growth-accounting decomposition,
\begin{equation}
  \Delta \log Y_{t} = \Delta \log A_{t} + \alpha \, \Delta \log K_{t}
    + (1 - \alpha) \, \Delta \log L_{t},
  \label{eq:growth}
\end{equation}
so measured output growth splits into a productivity term and the two factor
terms weighted by $\alpha$ and $1 - \alpha$.

The result worth stating is what~\eqref{eq:growth} leaves behind: whatever is not
explained by capital and labour is charged to $A_{t}$, which is why the residual
carries so much of the action in practice. This paragraph also lives in
\texttt{sections/results.tex}, giving a second file to click on.

\subsection{Point estimates}

Estimating~\eqref{eq:growth} on the panel of Section~\ref{sec:data} by
least squares with unit and year effects gives $\hat{\alpha} = 0.34$ with a
standard error of $0.02$, comfortably inside the range the calibration
literature treats as uncontroversial. The implied productivity residual accounts
for a little under half of average output growth, and rather more than half of
its variance across economies.

\begin{table}[ht]
  \centering
  \begin{tabular}{lrrr}
    \hline
    & \multicolumn{1}{c}{(1)} & \multicolumn{1}{c}{(2)} & \multicolumn{1}{c}{(3)} \\
    \hline
    $\Delta \log K$  & 0.34 & 0.31 & 0.36 \\
                     & (0.02) & (0.03) & (0.02) \\
    $\Delta \log L$  & 0.66 & 0.69 & 0.64 \\
                     & (0.02) & (0.03) & (0.02) \\
    Unit effects     & Yes & Yes & Yes \\
    Year effects     & No  & Yes & Yes \\
    Balanced sample  & No  & No  & Yes \\
    \hline
    Observations     & 1{,}512 & 1{,}512 & 1{,}280 \\
    $R^{2}$          & 0.61 & 0.68 & 0.71 \\
    \hline
  \end{tabular}
  \caption{Growth-accounting regressions. Standard errors clustered by unit in
    parentheses; factor shares constrained to sum to one.}
  \label{tab:estimates}
\end{table}

Table~\ref{tab:estimates} shows that the estimate is not delicate. Adding year
effects moves $\hat{\alpha}$ by three points, restricting attention to the
balanced sample moves it by two in the other direction, and no specification
takes it outside $[0.30, 0.37]$. Given the exponent in~\eqref{eq:kstar}, however,
even that narrow band maps into a range for $K^{\star}$ of roughly twenty per
cent, which is a good deal less comforting than the standard errors suggest.

\subsection{Transition dynamics}

Simulating the model forward from a capital stock half its steady-state value
produces the familiar concave path: output growth starts near six per cent and
decays geometrically toward zero. The half-life of the gap is
\begin{equation}
  h = \frac{\log 2}{\lambda},
  \qquad
  \lambda \approx (1 - \alpha)\left( \beta^{-1} - 1 + \delta \right),
  \label{eq:halflife}
\end{equation}
which at the estimated parameters is a little over eighteen years. Convergence
in this class of models is slow, and~\eqref{eq:halflife} makes clear that it is
slow for a structural reason rather than an accident of calibration: raising
$\alpha$ lowers $\lambda$ directly.

The gap between the model's transition and the data's is where an honest draft
starts hedging. The simulated paths are smoother than the observed ones, they
lack the persistence of the measured residual, and they cannot reproduce the
handful of large contractions that survive in the sample. A reader who wants to
argue with any of that will find this paragraph a convenient place to start.
